
\chapter{Introduction}

\section{The problem}

A photograph of a diseased leaf carries the information needed to name the disease.
Teaching a convolutional network to recover that information from labelled examples
is now routine, and published accuracies on curated datasets sit above 99\,\%. The
number is impressive and it is also fragile, because the datasets that produce it
were collected under narrow conditions: single leaves, plain backgrounds, even
lighting, one disease per frame.

A farmer photographs a leaf on a plant, in a field, with soil and other leaves
behind it and the sun somewhere off-axis. If a model trained on the first kind of
image is evaluated on the second, the accuracy it reports and the accuracy it
delivers are different quantities.

This project measures both, then asks why they differ. The distinction matters
because a screening tool that is wrong four times out of five is worse than no tool
at all, and because a model that appears to work is harder to question than one that
visibly fails.

\section{Motivation}

Plant disease causes substantial annual losses, and early identification is one of
the cheaper interventions available. Automated identification from a phone
photograph is attractive precisely where expert advice is scarce, which is also
where the consequences of a wrong answer are largest.

Three things make the gap between reported and delivered accuracy worth studying
rather than merely noting. First, it is large: our own measurement is a 76.8 point
drop in macro-F1 between curated and field images. Second, it is invisible in the
usual reporting: a paper that quotes only curated accuracy is not lying, it is
answering a different question. Third, an explanation exists in the literature and
is usually asserted rather than tested. If models are exploiting laboratory
backgrounds, targeted augmentation should help, and that prediction is testable.

\section{Objectives}

\begin{enumerate}[leftmargin=1.4em,itemsep=2pt]
\item Train a CNN classifier on PlantVillage under a protocol frozen before any
      result is inspected, and report held-out performance with class-wise detail.
\item Measure the same checkpoints on PlantDoc without retraining, and quantify the
      domain gap with confidence intervals rather than a single number.
\item Test, rather than assume, whether the gap is caused by the model relying on
      background context, using an intervention experiment that needs no human
      labelling of saliency maps.
\item Test whether targeted background augmentation closes the gap, with three
      seeds per condition and a pre-registered direction of effect.
\item Bound the usefulness of the system with calibration and abstention analysis,
      and state plainly where it should not be used.
\end{enumerate}

\section{Scope}

The study covers one controlled dataset (PlantVillage, 38 classes) and one field
dataset (PlantDoc, 236 test images, 27 classes scored after label review). Four
architectures are compared under one protocol. The unit of independence for
inference is the training seed, not the image.

Work outside scope is stated in Chapter~\ref{ch:limitations}: no expert agronomic
re-labelling, no multi-disease frames, no severity grading, no deployment claim, and
no assertion of field readiness.

\section{Contributions}

\begin{enumerate}[leftmargin=1.4em,itemsep=2pt]
\item A measured domain gap with intervals, replicated over five seeds, rather than
      a single accuracy pair.
\item An intervention-based mechanism result that contradicts the common
      background-shortcut explanation in its strong form, and supports a weaker one.
\item A pre-registered null result for background augmentation, with the mechanism
      measurement used to explain why the fix cannot work.
\item A calibration and abstention bound showing the confidence threshold does not
      rescue the field case.
\item A repository in which every reported number is generated from a stored
      artifact by a registry check, so the document cannot drift from the evidence.
\end{enumerate}

\section{Organisation of the report}

Chapter~\ref{ch:literature} reviews prior work and identifies what is usually
asserted without measurement. Chapter~\ref{ch:methodology} describes the frozen
protocol, the architectures and the evaluation design. Chapter~\ref{ch:datasets}
documents both datasets, the label mapping decisions and the split construction.
Chapter~\ref{ch:results} answers each hypothesis in turn. Chapter~\ref{ch:discussion}
interprets the combined picture. Chapter~\ref{ch:limitations} states what the work
does not establish, including two defects found in the original pipeline.
Chapter~\ref{ch:conclusion} closes with what would be worth doing next.
